Toen nog hartje zomer
verloor zich de dromer

in haar kalm verleden
een paar jaar geleden.

Ze groeide voor haar vlucht
op in een klein gehucht.

Daar werd door rust verwarmd
de tijdloosheid omarmd.

Daar bracht de zon een lach
naar wie in het veld lag.

De lucht was dan zo zacht
dat men slechts heilzaam dacht.

Zag men de zon rijzen,
vogels het licht prijzen,

via het zonnelied,
dat Franciscus hen liet,

door hun fluitend gezang.
Zo duurde een dag lang.

Zag men door boomtakken
de enkele wrakken

die een brink omringden.
Zij die slootjespringden

omdat de brug rot was,
wilden het plein van gras

voor hun plaggenhutten,
niet zomaar benutten.

In deze contreien
vormden “boerderijen”,

of wat er voor door ging,
de gangbare woning.

De muren van keien,
die men wist te kleien,

ondersteunden maar net
het dak boven hun bed.

De ruim bemeten hal
diende vaak als een stal.

Verderop lag de es
waar, lopend in een hes,

Emmy’s vader oogste.
Zijn veld lag het hoogste,

waardoor het zeer vruchtbaar
en daardoor zeer kostbaar

om te bezitten was.
Ze waren krap bij kas

dus werd het land gepacht.
Dat gaf de landheer macht,

zoals haar moeders pracht
voor de huwelijksnacht.

Haar beklonken moeder,
die pa zag als loeder

en als verplicht vee hield,
kende ze steeds geknield.

Oogstend, biechtend, boenend
of een zwengel zoenend.

Het sluwe herenrecht
gold evenwel niet echt

in die verlichte tijd.
Het was zijn nep beleid

gebaseerd op Emer.
Nessa als afnemer,

zo luidt de legende.
Al ging het geplande

in de sage niet door,
het vond wel zijn gehoor

bij de valse landheer.
De rechtspraak van weleer

bood geen merkbaar verzet.
Hij ging vrijuit naar bed

met al zijn horigen,
meestal nog jeugdigen.

De regio typeert
als vrij geïsoleerd,

dus kwam het uitbuiten
mondjesmaat naar buiten.
De omgekochte proost
bood eveneens geen troost.

Hij verzweeg het misbruik,
al gaf zijn onderbuik,

die altijd rammelde
als zijn hart spartelde,

een andere boodschap.
De steun van het graafschap

vlak nabij het bisdom
was voor hem nodig om

de unieke kans op
het ambt van prins-bisschop

bindend te vergroten.
Familiegenoten

van de stoute landheer
voerden daar het beheer.

Dus werd niets bekend
en zelfs glashard ontkend.

Het mocht slechts eenmalig.
Doch werd hij inhalig

naar sommige vrouwen
van wie hij ging houden.

Zo werd haar moederlief
de landheers hartedief.

Maar meer mocht hij niet doen.
Ze hadden rechten toen.

Zowel voor heer als knecht
gold het Canoniek recht.

De pacht vroeg veel inbreng.
De eisen waren streng

en het achilleshiel
als de oogst tegenviel.

Kon men niet betalen,
dan kwam men verhalen

met als stevige dwang
de landheer zijn gevang.

Zo liet men zich knechten
en verloor zijn rechten

op het eigen bedrijf
behalve van het lijf.

Deze schuldslavernij
vierde toen nog hoogtij.

De boer door weer verweerd
was listig gechanteerd

en gewend te schikken.
“Zwijgen en ja-knikken”.

Met vieren onder dak
nam hij elk klein gemak

dat beter had gekund
alsof het werd vergund.

Tot nooddruft zijn gedoemd,
werd met geloof verbloemd.

Villa Volta tolde
in die hersenholte,

als maalstroom van ootmoed,
achterdocht en rampspoed,

tot die zure veenvloek,
uit het leenrecht-wetboek,
door een hogere macht,
eindelijk werd ontkracht.

Dan werd de belijder
pas een bokkenrijder.

Aldus, achttien jaar lang,
liep Emmy haar kruisgang.

Ze hoorde steeds klagen
over onbehagen

zonder dat haar ouders
maar ergens hun schouders

onder durfden zetten
om zich te verzetten.

Eenmaal een dubbeltje
werd je nooit een kwartje.

Ze kende elke hoek
van de dramadriehoek,

die in haar familie,
zonder introspectie,

dagelijks werd gespeeld.
Elke kaart was gedeeld.

Satan als reclamant
deelde de eerste hand.

Christus als bevrijder
en de mens als lijder.

Moeder was slachtoffer.
Haar gereedschapskoffer

bestond uit beschamen.
Pa moest het beamen

en was de aanklager.
Haar zus werd de drager

van schaamte, schande, schuld
en engelengeduld.

Want moeders volhardde
als ze haar spel startte.

Emmy was de redder
en daarmee de shredder

van al het ongemak.
Tot de last haar opbrak.

Tot ze de harde muur
van de klemmende schuur,

die haar moeder bouwde,
niet meer blind vertrouwde.

Haar wrede opvoeders
waren trouwe broeders

in strenge lijfstraffen.
Nevens het afblaffen

werd een houten lepel
of een leren gordel

gebruikt als opvoeding
of als vernedering.

Doch “Heb uw ouders lief”
smoorde het ongerief.

Elk kind hield zijn snavel.
Dit rook sterk naar zwavel.

Een slaventheater
vond Nietzsche pas later.

Wie doorbrak die tendens
en werd een Übermensch?
